\documentclass[11pt]{article}

% Venue-neutral double-blind draft. Replace this preamble with the official
% ICLR 2027 style after the venue releases it; do not change the evidence text.
\usepackage[margin=1in]{geometry}
\usepackage{amsmath}
\usepackage{booktabs}
\usepackage{graphicx}
\usepackage{hyperref}
\usepackage{multirow}
\usepackage[numbers,sort&compress]{natbib}

\title{Contract-Gated Adaptation and Radeon Deployment of
$\pi_{0.5}$ for Parcel Manipulation}
\author{Anonymous Authors}

\begin{document}
\maketitle

\begin{abstract}
% Write this paragraph last. It must contain the frozen LIBERO comparison,
% the frozen pure-VLA parcel result, and the Radeon efficiency result. Leave
% it empty until all three claim gates pass.
\end{abstract}

\section{Introduction}

Vision-language-action (VLA) policies promise a common interface for semantic
instruction following and continuous robot control, but contact-rich deployment
exposes a gap between plausible action generation and complete task success.
Parcel handling is a particularly demanding instance because thin mailers,
upright cylinders, horizontal tubes, and heavy cartons require different contact
and load-sharing behavior. A policy that reaches a destination but fails to form
a seal, retain the parcel, or release it safely has not completed the task.

This work studies $\pi_{0.5}$~\citep{pi05} as the learned policy for a mobile
bimanual parcel robot and as an executable baseline on LIBERO~\citep{libero}. We
separate public benchmark reproduction, task-specific development, frozen parcel
confirmation, and Radeon runtime measurements. This separation prevents public
simulation success from being presented as parcel success and prevents expert or
deterministic planner behavior from being credited to the VLA.

The proposed system combines an explicit observation--action contract, structured
absolute dual-arm actions, stage-aware flow training, verified recovery replay,
and a fail-closed promotion ladder. Deterministic code may transform frames,
interpolate commands, enforce force and workspace limits, or reject unsafe
actions, but it may not query an expert pose, select a grasp family, or complete a
failed action. Each architectural and runtime change is evaluated as a single
factor under paired initial states.

The paper is organized around three questions. First, can a contract-gated
adaptation improve over the public $\pi_{0.5}$ LIBERO baseline without using
privileged state? Second, can the same methodology produce reliable pure-VLA
pick--transport--place--release behavior across a heterogeneous parcel set?
Third, which Radeon runtime choices improve latency or energy without regressing
closed-loop success? The current draft reports completed baseline and negative
results and leaves all unfinished claim cells as \texttt{TBD}.

\section{Related Work}

$\pi_{0.5}$ extends the $\pi_0$ family with heterogeneous pretraining and
open-world generalization~\citep{pi05}. OpenPI provides base and task-specific
checkpoints, including a LIBERO model and a DROID-derived checkpoint used for
robot adaptation~\citep{openpi}. Large collections such as Open
X-Embodiment~\citep{openx} and DROID~\citep{droid} motivate cross-embodiment
pretraining, while LIBERO provides 130 language-conditioned manipulation tasks
with fixed initial states for controlled evaluation~\citep{libero}.

Continuous action chunks improve control throughput but become stale when visual
or contact state changes. OpenVLA-OFT studies continuous action prediction and
multi-image observations~\citep{openvlaoft}, while real-time chunking treats
overlapping flow-policy chunks as an inpainting problem~\citep{rtc}. We evaluate
execution cadence separately from checkpoint training because a faster or more
reactive runtime is not an architectural improvement unless success is preserved.

Interactive learning methods address compounding policy errors by collecting
states induced by the learned policy~\citep{dagger,ript}. RECAP applies
advantage-conditioned learning to demonstrations, corrections, and autonomous
experience in the $\pi_{0.6}^{*}$ system~\citep{recap}. In the frozen software
stack used for our experiments, we did not identify public $\pi_{0.6}$ weights
with a compatible LIBERO processor and action adapter. We therefore use
$\pi_{0.6}^{*}$ to motivate verified recovery experiments rather than as a
numerical baseline.

\section{Method}

\subsection{Policy Ownership and Contracts}

Every result is bound to checkpoint, processor, training-contract, dataset,
code, evaluation-design, and ROCm runtime hashes. A rollout is classified as
pure VLA only when the learned policy supplies every task decision and action in
the declared contract, no expert action or target pose is consumed, no fallback
completes the task, and the force, destination, and release predicates all pass.
A deterministic safety rejection remains a policy failure; action replacement
is reported as assisted control and receives no pure-VLA credit.

The parcel observation contract contains synchronized overhead and wrist RGB-D,
robot state, tool seals, force history, parcel geometry estimates, task phase,
and retry index. Missing sensors are masked explicitly. The 23-dimensional
training target contains a three-dimensional base-velocity command, absolute
world-frame positions and unit quaternions for both end effectors, binary tool
commands, three learned grasp-mode logits, and primitive progress. The executable
decoder normalizes each quaternion, clips base velocity to predeclared limits,
and rejects nonfinite fields; round-trip tests cover normalization,
denormalization, decoding, bounds, and phase transitions.

\subsection{Stage-Aware Flow Adaptation}

The baseline predicts a length-$H$ continuous action chunk with a flow-matching
objective. Parcel demonstrations are imbalanced across pregrasp, grasp, lift,
transport, place, and terminal phases. We therefore apply a pre-registered stage
weight $w_s$ only to the per-sample continuous flow loss,
\begin{equation}
  \mathcal{L}=w_s\mathcal{L}_{\mathrm{flow}}+
  \lambda_{\mathrm{mode}}\mathcal{L}_{\mathrm{mode}},
\end{equation}
while leaving the categorical mode loss unweighted. The weights are normalized
so that their dataset-frequency-weighted mean equals one. This preserves the
overall gradient scale and isolates the effect of emphasizing contact-critical
stages from changes to the already trained mode decoder.

\subsection{Verified Recovery Replay}

Failed VLA actions are retained as failure telemetry rather than positive
behavior-cloning labels. A corrective trajectory enters the supervision pool
only after independent replay completes grasp, retention, transport, placement,
and release without target leakage, expert fallback, or safety violations. The
next model is promoted only when paired frozen development improves and safety,
energy, memory, and attribution gates do not regress. We refer to this procedure
as verified replay or gated self-improvement, not autonomous self-evolution,
unless the learned system generates and verifies its own corrections.

\subsection{Radeon Runtime Selection}

Inference uses a persistent policy process and separates checkpoint load, model
inference, queued action dispatch, and observation-to-actuation latency. Runtime
experiments record warm P50/P95/P99 latency, peak memory, package energy,
junction temperature, sampling errors, and complete task outcomes. Compilation
or a changed action-execution cadence is accepted only when paired success does
not regress and the pre-registered latency and energy gates both pass.

\section{Experimental Design}

\subsection{Public LIBERO Benchmark}

The public evaluation uses LIBERO-Spatial, LIBERO-Object, LIBERO-Goal, and
LIBERO-10 with ten tasks per suite and ten fixed initial states per task, giving
400 independent rollouts. The official $\pi_{0.5}$ checkpoint uses relative
seven-dimensional actions, 50-step predicted chunks, and ten executed actions per
inference. Initial states 10--19 form a development population; states 0--9 are
reserved for confirmation. Candidate and control policies use identical task
order, seeds, checkpoint, observation adapters, and success predicates.

The primary endpoint is complete sparse task success. Single-model proportions
are reported with Wilson 95\% intervals. Paired candidates are compared with an
exact one-sided McNemar test and a paired Newcombe interval. Published aggregate
scores provide context only; a superiority claim requires paired local outcomes
on the same 400 confirmation units.

\subsection{Frozen Parcel Evaluation}

The parcel confirmation design stratifies object family, size, mass, pose,
contact mode, destination, lighting, and perturbation. Complete success requires
acquisition, stable retention, transport, correct placement, release, force
compliance, and no expert fallback. Intermediate seal, lift, retention,
destination, release, and safety rates are diagnostic outcomes, while complete
pure-VLA success remains the sole primary endpoint. Development cases may guide
checkpoint selection; confirmation cases are opened once after the configuration
is frozen.

\subsection{Ablations}

The study changes one declared factor at a time. B2 uses the absolute action
contract, B3 replaces it with bounded incremental SE(3), B3-SW adds normalized
stage weighting to the continuous flow loss, B4 adds verified recovery replay,
and B5 adds short observation memory and explicit policy metadata. Runtime
ablations keep the checkpoint fixed and compare action-execution cadence,
persistent serving, and eager or compiled execution.

\section{Results}

\subsection{Public Baseline Reproduction}

The immutable local $\pi_{0.5}$ LIBERO checkpoint completed 385 of 400 frozen
confirmation rollouts (96.25\%). Suite success was 98\% on LIBERO-Spatial,
99\% on LIBERO-Object, 96\% on LIBERO-Goal, and 92\% on LIBERO-10. The
published OpenPI aggregate is 96.85\%~\citep{openpi}; it is not paired with the
local run and is therefore reported only as context. Table~\ref{tab:libero}
reserves the candidate row until development selection and one-shot confirmation
are complete.

\begin{table}[t]
\centering
\caption{LIBERO success under the four-suite protocol. Published results are
context only; local superiority requires paired confirmation outcomes.}
\label{tab:libero}
\begin{tabular}{@{}lccccc@{}}
\toprule
System & Spatial & Object & Goal & LIBERO-10 & Mean \\
\midrule
OpenPI $\pi_{0.5}$, published & 98.8 & 98.2 & 98.0 & 92.4 & 96.85 \\
Local $\pi_{0.5}$, 400 confirmation & 98.0 & 99.0 & 96.0 & 92.0 & 96.25 \\
Selected candidate, 400 confirmation & TBD & TBD & TBD & TBD & TBD \\
\bottomrule
\end{tabular}
\end{table}

\subsection{Action-Execution Cadence Screen}

The frozen cadence screen evaluated four executed action prefixes on 40 paired
development contexts. Eight executed actions per inference achieved 40/40 sparse
successes, compared with 38/40 for the ten-action control, and was therefore the
only arm selected for full development evaluation. This is a screening result,
not a confirmation score. Shortening the prefix to four actions reduced warm
model P95 by 6.3\% but did not improve success and nearly doubled package energy;
reactivity alone was therefore insufficient for promotion.

\begin{table}[t]
\centering
\caption{Frozen LIBERO cadence screen on 40 development contexts. Selection
maximizes complete success first and prefers a longer prefix only to break a
success tie. The selected eight-action arm must still pass 400 paired development
units before confirmation.}
\label{tab:action-steps}
\begin{tabular}{@{}lrrrr@{}}
\toprule
Executed actions & Success & Model calls & Warm P95 & Energy \\
\midrule
10 (control) & 38/40 & \textbf{634} & 528.58 ms & \textbf{18.37 Wh} \\
4 & 38/40 & 1,589 & \textbf{495.03 ms} & 35.70 Wh \\
\textbf{8 (selected)} & \textbf{40/40} & 751 & 632.84 ms & 20.72 Wh \\
6 & 39/40 & 1,040 & 536.59 ms & 25.33 Wh \\
\bottomrule
\end{tabular}
\end{table}

\subsection{Action-Contract Gate}

The 12k-step B3 incremental-SE(3) candidate did not pass the offline action gate.
On six paired observations it improved two and regressed four relative to the B2
absolute-action control. Its mean normalized pose-error improvement was $-0.546$
(negative values indicate regression), and the one-sided promotion checks failed.
No B3 parcel closed-loop score is reported. This result motivates B3-SW but does
not by itself establish that stage weighting will repair the action contract.

\subsection{Radeon Runtime and Endurance}

The eager 400-rollout LIBERO run used 8.97~GiB peak allocated memory, consumed
168.31~Wh including process startup, and had a warm model-inference P95 of
518.55~ms. A paired ABBA development experiment found that max-autotune
compilation reduced median warm P95 by 51.19\%, but it regressed paired success in
one comparison and worsened median energy per episode by 7.06\%. Compilation was
therefore rejected despite its latency improvement.

The selected eager runtime completed three cycles over the same 40 task/state
contexts with 37, 38, and 38 successes. These 120 attempts contain only 40
independent contexts and are reported as endurance evidence rather than a larger
benchmark sample. Median episode P95 drifted by 0.63\% from the first to third
cycle, peak junction temperature was 52~$^\circ$C, and steady-state energy was
0.347~Wh per pure-VLA success.

\begin{table}[t]
\centering
\caption{Verified Radeon evidence. Lower latency alone does not promote a runtime.}
\label{tab:radeon}
\begin{tabular}{@{}lccc@{}}
\toprule
Experiment & Success evidence & Warm model P95 & Energy \\
\midrule
Eager LIBERO confirmation & 385/400 & 518.55 ms & 168.31 Wh \\
Max-autotune ABBA & success gate failed & $-51.19\%$ & $+7.06\%$/episode \\
Eager endurance & 37/40, 38/40, 38/40 & 0.63\% drift & 0.347 Wh/success \\
\bottomrule
\end{tabular}
\end{table}

\subsection{Incomplete Results}

The full paired LIBERO action-step development comparison, one-shot candidate
confirmation, B3-SW action gate, frozen parcel confirmation, verified-replay
ablation, and final Radeon action-runtime comparison are \texttt{TBD}. They will
be inserted only from machine-readable summaries after their pre-registered
gates complete.

\section{Discussion}

The completed experiments show why VLA papers should separate checkpoint
integrity, offline action fidelity, closed-loop success, and accelerator speed.
The B3 checkpoint trained successfully and passed static audits yet failed its
first behavioral interface. Conversely, compilation produced a large latency
gain but was not deployable because success and energy gates failed. These
negative results are scientifically useful because they rule out two common but
unsupported shortcuts: treating training completion as robot competence and
treating kernel speed as system improvement.

The study currently has three important limitations. First, the public result is
a single-model LIBERO reproduction, whereas the parcel claim requires a separate
task-specific confirmation population. Second, the parcel dataset contains few
independent demonstrations relative to its frame count, so training changes must
be interpreted as adaptation under limited behavioral coverage. Third, we did not
identify public $\pi_{0.6}$ weights with a processor and evaluation adapter
matching the frozen stack; a numerical superiority claim would therefore be
unverifiable under this protocol. The final paper will compare directly with
executable $\pi_{0.5}$ baselines and discuss $\pi_{0.6}^{*}$ only as
methodological motivation.

\section{Conclusion}

% Write after the candidate LIBERO and frozen parcel gates complete.

\bibliographystyle{plainnat}
\bibliography{references}

\end{document}
